%% alg-claim.tex -- Algorithm: the claim, as the loops actually write it.
%%
%% An earlier version of this figure showed `FOR UPDATE SKIP LOCKED'. No loop in
%% the implementation uses it: every one selects a candidate and then issues a
%% conditional UPDATE whose WHERE clause still names the old status. That predicate
%% is what does the arbitration, and it is worth showing literally, because the
%% property it establishes is narrower than the one the earlier draft claimed.
%%
%% The property is concurrent mutual exclusion, not exactly-once execution across
%% crashes. Table tab:loops states what each loop does after a crash, which is where
%% the two come apart.

\begin{algorithm}[tbp]
\caption{The claim, as the loops write it. The arbitration is the conditional
\code{UPDATE} on lines~4--7: two workers may both select row $r$, but only one
\code{UPDATE} finds it still \plit{queued}, and the loser sees
\code{rowcount}~$=0$. What this establishes is \emph{concurrent mutual exclusion} ---
no two workers own a row at the same instant. It does not establish exactly-once
execution across a crash, because what happens to a row whose owner dies is a
per-loop recovery policy, not a property of the claim (\tabref{tab:loops}).}
\label{alg:claim}
\psig{claim}
\pin{$Q$}{a queue table with a \pid{status} column and, on some loops,
  \pid{claimed\_by} and \pid{last\_heartbeat\_at}}
\pin{$w$}{this worker's identity}
\pinv{no two workers hold the same row in \plit{running} at the same instant, for
  any number of replicas}
\begin{pseudo}
\pline{$r \gets$ \plit{SELECT} \kw{one} \kw{from} $Q$ \kw{where}
       \pid{status} $=$ \plit{queued} \kw{order by} \pid{enqueued\_at}}
\pline{\kw{if} $r = \bot$ \pthen{} \kw{wait}; \kw{return} $\bot$}
\pline{}
\pline{$n \gets$ \plit{UPDATE} $Q$ \plit{SET} \pid{status} $=$ \plit{running},
       \pid{claimed\_by} $= w$}
\pline{\I \plit{WHERE} \pid{id} $= r.\pid{id}$
       \kw{and} \pid{status} $=$ \plit{queued}
       \cmt{the arbitration}}
\pline{\kw{commit}}
\pline{\kw{if} $n = 0$ \pthen{} \kw{return} $\bot$
       \cmt{another worker won it}}
\pline{}
\pline{\kw{for each} \pid{step} \kw{in} stages($r$) \pdo}
\pline{\I run(\pid{step}); persist(\pid{step}); \kw{commit}}
\pline{\I \kw{if} $Q$ \kw{has} \pid{last\_heartbeat\_at} \pthen{} touch it}
\pline{\pid{status} $\gets$ \plit{done}; \kw{commit}}
\end{pseudo}

\vspace{0.6ex}
\pinv{After the owner of a running row dies, the row's fate is decided by that
  loop's recovery policy. The refinement janitor marks it \plit{failed}; the
  workflow-run worker requeues it on lease expiry; the calibration drain records no
  lease and so has no automatic recovery at all. \tabref{tab:loops} is the
  authority, and the differences are why this paper no longer states a single
  delivery guarantee for ``the queue''.}
\end{algorithm}
